\documentclass[12pt,letterpaper]{article}
\usepackage[utf8]{inputenc}

% Ajustamos los márgenes para corregir las advertencias de fancyhdr
\usepackage[letterpaper, margin=1in, bottom=3in, top=2.2in]{geometry}

% Cargamos primero el paquete fancyhdr
\usepackage{fancyhdr}
% Corregimos los parámetros headheight y footskip
\setlength{\headheight}{15pt} % Al menos 14.49998pt según la advertencia
\setlength{\footskip}{72pt}   % Al menos 71.60547pt según la advertencia

\usepackage{xcolor}
\usepackage{amsmath}
\usepackage{anyfontsize}
\usepackage{graphicx}
\usepackage{tikz} % Necesario para los elementos gráficos
\usepackage{listings} % Para resaltado de sintaxis
\usepackage{sourcesanspro}
\renewcommand{\familydefault}{\sfdefault}
\usepackage{fontawesome5}
\usepackage{titlesec}
\usepackage{setspace}
\usepackage{hyperref}
\usepackage{mdframed} % Para marcos más estables
\usepackage{qrcode} % Paquete para generar códigos QR
\usepackage{eso-pic}

% Cargar tikz y la biblioteca de cajitas para mdframed
\usetikzlibrary{calc,shapes,positioning}

% IMPORTANT: Force the text color to be white for the entire document
\AtBeginDocument{\color{primaryColor}}

% Colors updated with better contrast
\definecolor{myblue}{cmyk}{1.0,0.8,0.05,0.20}
\definecolor{bgColor}{RGB}{15, 22, 36}  % Very dark blue, almost black background
\definecolor{primaryColor}{RGB}{255, 255, 255}  % White text - full white for better contrast
\definecolor{accentColor}{RGB}{255, 212, 59}  % Python Yellow
\definecolor{pythonBlue}{RGB}{120, 180, 255}  % Azul Python más claro para mejor contraste
\definecolor{secondaryColor}{RGB}{220, 230, 240}  % Lighter blue-gray for better contrast
\definecolor{terminalBg}{RGB}{22, 22, 22}  % Terminal background remains dark
\definecolor{terminalFrame}{RGB}{40, 40, 40}  % Terminal frame slightly darker
\definecolor{lineNumberColor}{RGB}{100, 100, 100}  % Color más sutil para los números de línea
\definecolor{dividerColor}{RGB}{60, 70, 90}  % Color más sutil para la línea divisoria

% Colores mejorados para el código con mayor contraste
\definecolor{codeTextColor}{RGB}{255, 255, 255}  % Blanco puro para el texto básico
\definecolor{codeKeywordColor}{RGB}{135, 206, 250}  % Azul cielo para palabras clave
\definecolor{codeCommentColor}{RGB}{144, 238, 144}  % Verde claro para comentarios
\definecolor{codeStringColor}{RGB}{255, 215, 135}  % Amarillo para strings
\definecolor{codeMethodColor}{RGB}{255, 160, 122}  % Naranja claro para métodos
\definecolor{codeFunctionColor}{RGB}{255, 215, 0}  % Dorado para funciones
\definecolor{codeNumberColor}{RGB}{255, 255, 150}  % Amarillo claro para números

% Set page background color
\pagecolor{bgColor}

% Hyperref configuration
\hypersetup{
    colorlinks=true,
    linkcolor=accentColor,
    filecolor=accentColor,
    urlcolor=accentColor,
}

% Configuración básica para listings que evita problemas con saltos de página
\lstset{
  language=Python,
  basicstyle=\ttfamily\bfseries\color{codeTextColor}, % Tamaño reducido, en blanco y negrita
  backgroundcolor=\color{terminalBg},
  commentstyle=\color{codeCommentColor},
  keywordstyle=\color{codeKeywordColor},
  stringstyle=\color{codeStringColor},
  numberstyle=\color{lineNumberColor}, % Números de línea más sutiles
  breaklines=true,
  breakatwhitespace=true,
  tabsize=4,
  showstringspaces=false,
  frame=none,
  xleftmargin=15pt,
  xrightmargin=0pt,
  aboveskip=10pt,
  belowskip=10pt,
  numbers=left,
  numbersep=8pt,
  extendedchars=true,
  keepspaces=true,
  columns=flexible,
  lineskip=6pt, % Espacio entre líneas ligeramente reducido
  % Keywords de Python (palabras reservadas)
  morekeywords={self, yield, lambda, with, as, from, True, False, None, import, in, for, 
                if, elif, else, while, return, def, class, try, except, finally, raise, 
                break, continue, global, nonlocal, pass, assert, del, yield from},
  % Funciones incorporadas resaltadas de manera distinta
  emph={[2]print,range,sum,int,str,float,list,dict,set,tuple,next,len,type,map,filter,
         reduce,zip,enumerate,sorted,reversed,min,max,open,any,all},
  emphstyle={[2]\color{codeFunctionColor}}
}

% TERMINAL ESTILO MAC MEJORADA CON BOTONES A LA IZQUIERDA Y MENOS ESPACIO SUPERIOR
\newenvironment{macterminal}{%
    \begin{mdframed}[
        linecolor=terminalFrame,
        backgroundcolor=terminalBg,
        roundcorner=5pt,
        skipabove=10pt,
        skipbelow=10pt,
        linewidth=1pt,
        innertopmargin=10pt, % Reducido
        frametitle={%
            \tikz[baseline=(current bounding box.east), outer sep=0pt]{
                \fill[red!80!black] (0,0) circle (5pt);
                \fill[yellow!80!black] (0.7,0) circle (5pt);
                \fill[green!70!black] (1.4,0) circle (5pt);
            }
        },
        frametitlealignment=\raggedright, % Alineado a la izquierda
        frametitleaboveskip=8pt, % Espacio entre el título y el borde superior
        frametitlebelowskip=0pt, % Reducir el espacio entre el título y el contenido
    ]
    % No necesitamos más espacio vertical aquí
}{%
    \end{mdframed}%
}

% Utilities
\newcommand{\verspace}{\vspace{10pt}}

% Section styling - Esquema más elegante y mejorado contraste
\titleformat{\section}
  {\LARGE\bfseries\color{primaryColor}} % Secciones principales en blanco para mejor contraste y más grandes
  {\thesection. }
  {0pt}
  {}
  []

\titleformat{\subsection}
  {\Large\bfseries\color{accentColor}} % Subsecciones en amarillo Python y más grandes
  {\thesubsection. }
  {0pt}
  {}
  []

\titleformat{\subsubsection}
  {\large\bfseries\color{pythonBlue}} % Subsubsecciones en azul Python más claro y más grandes
  {\thesubsubsection. }
  {0pt}
  {}
  []

% Spacing for sections
\titlespacing*{\section}{0pt}{20pt}{10pt}
\titlespacing*{\subsection}{0pt}{15pt}{7pt}
\titlespacing*{\subsubsection}{0pt}{10pt}{5pt}

% Configure header and footer
\pagestyle{fancy}
\fancyhf{}

% Crear un encabezado elegante para todas las páginas (excepto la portada)
\fancyhead[L]{
    \begin{tikzpicture}[remember picture, overlay]
        % Rectángulo de fondo para la sección Python
        \fill[pythonBlue, rounded corners=3pt] (0,0) rectangle (2.2cm,0.7cm);
        % Texto Python
        \node[text=primaryColor, font=\bfseries] at (1.1cm,0.35cm) {PYTHON};
        
        % Pequeño separador
        \fill[secondaryColor] (2.35cm,0.1cm) rectangle (2.40cm,0.6cm);
        
        % Texto Concurrency & Parallelism con el color de acento
        \node[text=accentColor, font=\bfseries, anchor=west] at (2.35cm,0.35cm) {CONCURRENCY (PART 2)};
    \end{tikzpicture}
}

% Línea de encabezado
\renewcommand{\headrulewidth}{0pt}


% Configure footer with profile info
\fancyfoot[C]{
    \vspace*{0.2cm}
    \noindent
    \begin{minipage}{\textwidth}
        \begin{flushleft}
            % Profile image placeholder with TikZ
            \raisebox{0.7cm}{
            \begin{tikzpicture}[baseline]
                \path[fill=bgColor] (0,0) circle (0.8cm);
                \clip (0,0) circle (0.8cm);
                \node at (0,0) {
                    % Placeholder for profile image
                    \includegraphics[width=1.6cm,height=1.6cm]{slides/profile-image.jpeg}
                };
            \end{tikzpicture}
            }
            \begin{minipage}[b]{0.8\textwidth}
                % Name
                {\large\bfseries\color{primaryColor}Alejandro Sánchez Yalí}
                
                % Professional description
                \par\vspace{1pt}
                {\small\color{secondaryColor}Software Developer | AI \& Blockchain Enthusiast}
                
                % Contact
                \par\vspace{1pt}
                {\small\color{accentColor}\faGlobe\hspace{5pt}\color{secondaryColor}www.asanchezyali.com}
            \end{minipage}
        \end{flushleft}
        \vspace{8pt} % Space at bottom
    \end{minipage}
}

% Estilo para la portada (sin encabezado ni pie de página)
\fancypagestyle{plain}{
    \fancyhf{}
    \renewcommand{\headrulewidth}{0pt}
    \renewcommand{\footrulewidth}{0pt}
}

% Bullet styling for lists
\renewcommand{\labelitemi}{\textcolor{accentColor}{$\bullet$}} % Primer nivel de viñetas en amarillo (mejor contraste)
\renewcommand{\labelitemii}{\textcolor{pythonBlue}{$\circ$}} % Segundo nivel de viñetas en azul

% Aumentar globalmente el tamaño de texto un 20%
\usepackage{relsize}
\AtBeginDocument{\relsize{1}} % Incrementa la fuente en aproximadamente un 20%

% Comando simplificado para la etiqueta Python
\newcommand{\languagetag}[1]{
    \begin{tikzpicture}[baseline]
        \node[fill=pythonBlue, text=primaryColor, rounded corners=5pt, inner sep=7pt] {
            {\normalsize\textbf{#1}}
        };
    \end{tikzpicture}
}

% Luego, define un comando para colocar la imagen de fondo en la primera página
% Coloca esto antes de \begin{document} (cerca de la línea 310-325)
\newcommand{\BackgroundPic}{%
    \put(0,0){%
        \begin{tikzpicture}[remember picture, overlay]
            % Coloca la imagen de fondo que cubra toda la página
            \node[inner sep=0pt] at (current page.center) {
                \includegraphics[width=\paperwidth,height=\paperheight,keepaspectratio=false]{slides/Coder_at_Night_in_the_City.jpeg}
            };
            
            % Agrega una capa semitransparente oscura sobre la imagen
            % Puedes ajustar el color y la opacidad según necesites
            \fill[black, opacity=0.7] (current page.south west) rectangle (current page.north east);
            
            % Si quieres un degradado oscuro en lugar de un color sólido, usa esto:
            % \shade[top color=black!80, bottom color=black!40, opacity=0.7] 
            %    (current page.south west) rectangle (current page.north east);
        \end{tikzpicture}%
    }%
}


% Alternativa usando una versión modificada de la URL
\newcommand{\elegantqr}[2]{
    \qrcode[height=2.5cm]{#1}
    \\[0.1cm]
    {\hspace{0.2cm}\color{primaryColor}\small #2\par}
}

% Definir el título personalizado con alineación exacta para todos los elementos
\newcommand{\titlepagecontents}{%
    \AddToShipoutPicture*{\BackgroundPic}
    \vspace*{2cm}
    % Asegurar que todos los elementos estén alineados a la izquierda
    \begin{flushleft}
    \languagetag{Python}\\[0.4cm]
    {\fontsize{48}{52}\bfseries\color{primaryColor}Python \color{accentColor}Concurrency\par}
    \vspace{0.3cm}
    {\fontsize{18}{52}\color{secondaryColor}Part 2: Asyncio and HTTP Operations\par}
    \vspace{0.3cm}
    {\color{secondaryColor}\today\par}
    \vspace{2cm}
    % El QR se alineará a la izquierda aunque internamente esté centrado
    \elegantqr{https://github.com/asanchezyali/social-media-posts/tree/main/Python/Concurrency}{Source Code}
    \end{flushleft}
}


% En la página del título, aplicar el contenido sin márgenes adicionales
\makeatletter
\renewcommand{\maketitle}{%
    \begin{titlepage}
    \setlength{\parindent}{0pt}
    \setlength{\leftskip}{0pt}
    \titlepagecontents
    \end{titlepage}
}
\makeatother

\newcommand{\finalpagecontents}{%
    \vspace*{4cm}
    \begin{center}
        \begin{tikzpicture}
            % Title text
            \node[text width=14cm, align=center] at (0,0) {
                {\fontsize{48}{52}\bfseries\color{primaryColor}Don't miss\\\color{accentColor}Part 3\par}
            };
            
            % Add vertical space
            \node at (0,-3) {};
            
            % Question text
            \node[text width=14cm, align=center] at (0,-5) {
                {\fontsize{24}{30}\color{secondaryColor}In the final part of our series, we'll explore\\\vspace{0.5cm}
                \color{accentColor}advanced patterns and hybrid approaches\\\vspace{0.4cm}
                \color{primaryColor}to take your Python concurrency skills to the next level!\par}
                
            };
        \end{tikzpicture}
    \end{center}
}

\begin{document}
\setlength{\parindent}{0pt}
\color{primaryColor} % Explicitly set text color to white for the entire document

\begin{titlepage}
    % Eliminamos todos los espacios adicionales para asegurar la alineación
    \titlepagecontents
\end{titlepage}

\section{Asyncio for Modern Concurrency}

\subsection{Coroutines and the Event Loop}

Asyncio provides a different concurrency model for I/O-bound tasks:

\begin{macterminal}
\begin{lstlisting}
import asyncio
import time

async def say_after(delay, message):
    """Coroutine that waits and then prints a message"""
    await asyncio.sleep(delay)
    print(message)

async def main():
    # Sequential execution
    print("Sequential execution:")
    start = time.time()
    
    await say_after(1, "Hello")
    await say_after(2, "World")
    
    print(f"Sequential took {time.time() - start:.2f} seconds\n")
    
    # Concurrent execution
    print("Concurrent execution:")
    start = time.time()
    
    # Create tasks to run concurrently
    hello_task = asyncio.create_task(say_after(1, "Hello"))
    world_task = asyncio.create_task(say_after(2, "World"))
    
    # Wait for both tasks to complete
    await hello_task
    await world_task
    
    print(f"Concurrent took {time.time() - start:.2f} seconds")

# Run the main coroutine
asyncio.run(main())

# Output
# Sequential execution:
# Hello
# World
# Sequential took 3.00 seconds

# Concurrent execution:
# Hello
# World
# Concurrent took 2.00 seconds
\end{lstlisting}
\end{macterminal}

This example demonstrates the basic concepts of asyncio:

\begin{itemize}
    \item \textbf{\textcolor{pythonBlue}{Coroutines:}} Functions defined with \texttt{async def} that can be paused and resumed
    \item \textbf{\textcolor{pythonBlue}{Tasks:}} Scheduled coroutines managed by the event loop
    \item \textbf{\textcolor{pythonBlue}{Await:}} The keyword that yields control back to the event loop until an operation completes
    \item \textbf{\textcolor{pythonBlue}{Event Loop:}} The central execution mechanism that coordinates all asyncio operations
\end{itemize}

The first example shows the difference between sequential execution (where operations block each other) and concurrent execution (where operations can run concurrently while one is waiting).

\subsection{Async with HTTP Requests}

A practical example of asyncio for concurrent HTTP requests:

\begin{macterminal}
\begin{lstlisting}
import asyncio
import aiohttp
import time

async def fetch_url(session, url):
    """Fetch a URL asynchronously"""
    start = time.time()
    async with session.get(url) as response:
        data = await response.text()
        elapsed = time.time() - start
        return {
            'url': url,
            'status': response.status,
            'size': len(data),
            'time': elapsed
        }

async def fetch_all(urls):
    """Fetch multiple URLs concurrently"""
    async with aiohttp.ClientSession() as session:
        tasks = [fetch_url(session, url) for url in urls]
        results = await asyncio.gather(*tasks)
        return results

async def main():
    # List of URLs to fetch
    urls = [
        'https://www.python.org',
        'https://docs.python.org',
        'https://www.github.com',
        'https://www.google.com',
        'https://www.wikipedia.org',
    ]
    
    # Sequential fetching for comparison
    start = time.time()
    async with aiohttp.ClientSession() as session:
        results = []
        for url in urls:
            results.append(await fetch_url(session, url))
    sequential_time = time.time() - start
    
    print(f"Sequential fetching: {sequential_time:.2f} seconds")
    
    # Concurrent fetching
    start = time.time()
    results = await fetch_all(urls)
    concurrent_time = time.time() - start
    
    print(f"Concurrent fetching: {concurrent_time:.2f} seconds")
    print(f"Speedup: {sequential_time/concurrent_time:.2f}x")
    
    # Display results
    for result in results:
        print(f"{result['url']}: {result['status']} ({result['size']} bytes, {result['time']:.2f}s)")

# Run the main coroutine
if __name__ == "__main__":
    asyncio.run(main())

# Output
# Sequential fetching: 2.26 seconds
# Concurrent fetching: 0.80 seconds
# Speedup: 2.83x
# https://www.python.org: 200 (51012 bytes, 0.12s)
# https://docs.python.org: 200 (17129 bytes, 0.11s)
# https://www.github.com: 200 (281297 bytes, 0.80s)
# https://www.google.com: 200 (16957 bytes, 0.46s)
# https://www.wikipedia.org: 200 (88783 bytes, 0.36s)
\end{lstlisting}
\end{macterminal}

This example shows how to use asyncio for a common real-world scenario: fetching multiple web pages concurrently. Key points:

\begin{itemize}
    \item \textbf{\textcolor{pythonBlue}{aiohttp:}} An async-compatible HTTP client/server for Python
    \item \textbf{\textcolor{pythonBlue}{asyncio.gather():}} Executes multiple awaitables concurrently and returns their results
    \item \textbf{\textcolor{pythonBlue}{Resource Management:}} Using \texttt{async with} to properly manage session resources
    \item \textbf{\textcolor{pythonBlue}{Performance Comparison:}} Direct comparison of sequential vs. concurrent approaches
\end{itemize}

This pattern is ideal for I/O-bound applications like web scrapers, API clients, and data fetchers.

\subsection{Producer-Consumer Pattern with Asyncio}

A classic concurrency pattern implemented with asyncio:

\begin{macterminal}
\begin{lstlisting}
import asyncio
import random

async def producer_consumer_example():
    # Create a queue with a maximum size
    queue = asyncio.Queue(maxsize=5)
    
    async def producer():
        """Produces tasks and adds them to the queue"""
        for i in range(10):
            # Create a task
            task = f"Task-{i}"
            
            # Put the task in the queue (will wait if queue is full)
            await queue.put(task)
            print(f"Producer: Created {task}")
            
            # Simulate variable production time
            await asyncio.sleep(random.uniform(0.1, 0.5))
        
        # Signal end of production
        print("Producer: Done producing")
    
    async def consumer(name):
        """Consumes tasks from the queue"""
        while True:
            try:
                # Get a task from the queue (will wait if queue is empty)
                task = await queue.get()
                
                # Check for sentinel value signaling end
                if task is None:
                    print(f"{name}: Received shutdown signal")
                    queue.task_done()
                    break
                
                print(f"{name}: Processing {task}")
                
                # Simulate processing time
                await asyncio.sleep(random.uniform(0.2, 0.7))
                
                print(f"{name}: Completed {task}")
                
                # Signal task completion
                queue.task_done()
                
            except asyncio.CancelledError:
                break
    
    # Start producer and consumer tasks
    producer_task = asyncio.create_task(producer())
    consumer_tasks = [
        asyncio.create_task(consumer(f"Consumer-{i}"))
        for i in range(3)
    ]
    
    # Wait for the producer to finish
    await producer_task
    
    # Wait for all items to be processed
    await queue.join()
    
    # Send shutdown signal to all consumers
    for _ in range(len(consumer_tasks)):
        await queue.put(None)
    
    # Wait for consumers to process shutdown signal
    await asyncio.gather(*consumer_tasks)

# Run the example
if __name__ == "__main__":
    asyncio.run(producer_consumer_example())
\end{lstlisting}
\end{macterminal}

The producer-consumer pattern is a fundamental concurrency pattern that asyncio implements with these key components:

\begin{itemize}
    \item \textbf{\textcolor{pythonBlue}{asyncio.Queue:}} A thread-safe queue designed for use in asyncio applications
    \item \textbf{\textcolor{pythonBlue}{Backpressure:}} Managing production rate with a maximum queue size
    \item \textbf{\textcolor{pythonBlue}{Consumer Pool:}} Multiple consumers processing items concurrently
    \item \textbf{\textcolor{pythonBlue}{Graceful Shutdown:}} Proper shutdown signals and waiting for all tasks to complete
\end{itemize}

This pattern is useful for scenarios like job queues, data processing pipelines, and message processing systems.

\section{Conclusion (Part 2)}

In this second part of our exploration of concurrency and parallelism in Python, we have focused on asyncio, a modern approach to concurrency:

\begin{itemize}
    \item We learned how asyncio provides single-threaded concurrency through its event loop architecture and coroutines.
    
    \item We explored practical applications of asyncio with HTTP requests using the aiohttp library, demonstrating significant performance improvements for I/O-bound operations.
    
    \item We implemented the classic producer-consumer pattern using asyncio.Queue, showing how to build efficient and well-managed data processing pipelines.
\end{itemize}

Asyncio represents a modern approach to concurrency in Python that is particularly well-suited for I/O-bound operations. Its elegant async/await syntax provides a readable and maintainable way to write concurrent code without the complexities of traditional threading approaches.

In Part 3, we will explore hybrid approaches that combine different concurrency models, advanced patterns, and performance optimization techniques to further enhance your Python concurrency skills.

\section{References (Part 2)}
\begin{itemize}
    \item Python Documentation. \textit{asyncio - Asynchronous I/O}. \href{https://docs.python.org/3/library/asyncio.html}{Link}
    
    \item Real Python. \textit{Async IO in Python: A Complete Walkthrough}. \href{https://realpython.com/async-io-python/}{Link}
    
    \item PyPI. \textit{aiohttp - Asynchronous HTTP Client/Server}. \href{https://pypi.org/project/aiohttp/}{Link}
    
    \item Lynn Root. \textit{Asyncio: We Did It Wrong}. \href{https://www.roguelynn.com/words/asyncio-we-did-it-wrong/}{Link}

    \item Luciano Ramalho. \textit{Fluent Python}, Chapter 18: Concurrency with asyncio. O'Reilly Media.
    
    \item Translated, Edited and written in collaboration with AI.
\end{itemize}

\section{Explore My Other Posts}
\vspace{10pt}
\begin{tikzpicture}
  % Section background - increased height slightly to ensure all content fits
  \fill[color=terminalBg, rounded corners=5pt] (0,0) rectangle (\textwidth,-5);
  
  % Calculate width available for text (total width minus QR code width minus margins)
  \pgfmathsetmacro{\availableWidth}{\textwidth-4cm}
  
  % Left column content
  \begin{scope}
    % Section title
    \node[text=accentColor, font=\Large\bfseries, anchor=north west] at (0.5cm,-0.5cm) {Enjoyed This Content?};
    
    % Subtitle
    \node[text=primaryColor, font=\normalsize, anchor=north west] at (0.5cm,-1.2cm) {Don't miss the first part of this series:};
    
    % Python post title - with highlighted style and limited width
    \node[text=secondaryColor, font=\large\bfseries, anchor=north west, text width=\availableWidth] at (0.5cm,-1.9cm) {\color{accentColor}Python Concurrency Part 1: \color{primaryColor}Understanding the GIL, Threading and Multiprocessing};
    
    % Brief description of the post - with limited width and adjusted position
    \node[text=secondaryColor, font=\normalsize, text width=\availableWidth, anchor=north west] at (0.5cm,-3.1cm) {
      Learn the fundamentals of concurrent programming in Python, how the Global Interpreter Lock works, 
      and when to use threading versus multiprocessing for optimal performance.
    };
  \end{scope}
  
  % QR Code on the right - vertically centered
  \node[anchor=center] at ({\textwidth-1.5cm}, {-2.5cm}) {
    \qrcode[height=2.2cm]{https://www.linkedin.com/feed/update/urn:li:activity:7310460714710691841/}
  };
  
\end{tikzpicture}
\vspace{5pt}

\clearpage
\thispagestyle{empty}
\finalpagecontents

\end{document}